\clearpage
\fancyhead[L]{\small CARLINI. EFFEMERIDI ASTRONOMICHE.}
\hypertarget{native-579-lower}{}
\begin{center}\small Göttingische gelehrte Anzeigen. Stück 10. Seiten 89..93. 19. Januar 1815.\\[3pt]\rule{.72\linewidth}{.3pt}\end{center}

{\setlength{\rightskip}{0pt plus 3em}\emph{Effemeridi astronomiche di Milano per l'anno} 1815 \emph{calcolate da} \textsc{Francesco Carlini}.
\emph{Con appendice. Dalla r.\,c. stamperia di governo.} Mailand 1814. Der Kalender umfasst 128 Seiten, der Anhang 118, im Oktavformat.\par}

Die regelmäßigen Beiträge zu dem astronomischen Kalender, der seine mustergültige Gliederung beibehalten hat, haben wir bei den früheren Jahrgängen besprochen. Diesmal sind sie um einen sehr wertvollen Beitrag ergänzt. Es ist das Verzeichnis der 34 \textsc{Maskelyne}-Sterne aus \textsc{Piazzi}s neuem Katalog, mit der für 1800 und 1850 berechneten Präzession, der Eigenbewegung und den Konstanten, die man zur Berechnung der Aberration und Nutation nach der Methode von Herrn \textsc{von Zach} braucht. So hat man alles beisammen, was zur Bestimmung des scheinbaren Orts für das ganze gegenwärtige und selbst für das vorhergehende Jahrhundert nötig ist. Der Anhang enthält die folgenden wichtigen Abhandlungen. Breite der Sternwarte Brera aus Beobachtungen der Zirkumpolarsterne von \textsc{Barnabas Oriani}. Nach einem gut geschriebenen Überblick über die Beobachtungsverfahren, die in der neueren praktischen Astronomie bis heute nach und nach zur möglichst genauen Bestimmung der Polhöhen eingesetzt wurden, wendet sich der hervorragende Astronom insbesondere der Mailänder Sternwarte zu. Er schildert die verschiedenen Versuche, ihre Polhöhe immer genauer zu bestimmen. Selbst die neuesten Beobachtungen, die \textsc{Carlini} mit einem 16 Zoll großen Repetitionskreis von \textsc{Lenoir} angestellt hatte — demselben, den \textsc{Mechain} in Barcellona und Montjoui verwendet hatte —, ließen noch eine Unsicherheit von einigen Bogensekunden zurück. Seit Ende 1810 besitzt die Sternwarte einen 3 Fuß großen Kreis mit stehender Säule von \textsc{Reichenbach}. Mit ihm wurde im Lauf eines Jahres eine ununterbrochene Reihe von Beobachtungen angestellt, die in den Jahrgängen 1812 und 1813 der Mailänder Ephemeriden vollständig abgedruckt sind. Nun erhalten wir die erste Ausbeute aus diesem reichen Schatz, nämlich die vollständigen Ergebnisse der Beobachtungen des Polarsterns, von $\delta$ in der Kassiopeia und von $\varepsilon$ im Großen Bären. Die Übersicht über die einzelnen Ergebnisse zeigt von Tag zu Tag noch kleine Unterschiede, die mehrere Bogensekunden erreichen. Sie dürften zum Teil auf unvermeidliche Beobachtungsfehler und zum Teil auf Veränderungen in der Atmosphäre zurückgehen, die Barometer und Thermometer nicht anzeigen, sofern man einen Teil der Unterschiede nicht der Bauart des Instruments selbst zuschreiben will. Die Art, wie der Kreis mit der Säule verbunden ist, hat bekanntlich einige Bedenken hervorgerufen. Deshalb wird \textsc{Oriani} noch eine bewegliche Libelle am Kreis selbst anbringen lassen, und es wird höchst interessant sein, die Wirkung dieser Änderung zu sehen. Der Polarstern wurde an 159 Tagen unterhalb des Pols und an 143 Tagen oberhalb des Pols beobachtet. Bemerkenswert ist, dass die Beobachtungen am Tag schlechter miteinander übereinstimmen und zugleich eine kleinere Zenitdistanz ergeben als die Beobachtungen bei Nacht. \textsc{Oriani} hat die Tagbeobachtungen daher ganz ausgeschlossen und aus den übrigen Folgendes abgeleitet
\SourceLeaf{580}
\begin{center}\small\begin{tabular}{@{}l@{\hspace{9pt}}r@{}}
Polhöhe der Sternwarte Brera $.\;.\;.\;.$ & $49^\circ\,28'\,0''713$ \\
Deklination des Polarsterns für 1811 $.\;.$ & $88\phantom{\circ}\,17\phantom{'}\,59.494$ \\
\end{tabular}\end{center}
\noindent Bildete man das Mittel aus \emph{allen} Beobachtungen, so ergäbe sich nach unserer Rechnung
\begin{center}\small\begin{tabular}{@{}l@{\hspace{9pt}}r@{}}
Polhöhe der Sternwarte $.\;.\;.\;.\;.\;.\;.$ & $49^\circ\,28'\,1''242$ \\
Deklination für 1811 $.\;.\;.\;.\;.\;.\;.\;.\;.$ & $88\phantom{\circ}\,17\phantom{'}\,59.583$ \\
\end{tabular}\end{center}
\noindent \textsc{Oriani} hat dieses Verfahren gewählt, weil er die Unterschiede bei den Tagbeobachtungen besonderen Veränderungen der Atmosphäre zuschreibt. Ein Teil davon mag aber jedenfalls auf die jährliche Parallaxe zurückgehen. Könnten wir äußere Einflüsse außer Acht lassen, die am Tag immer oder wenigstens überwiegend in derselben Richtung wirken, so würden wir dazu neigen, dem Ergebnis unserer Rechnung, einer maximalen Parallaxe in Deklination von $0''7$, einiges Gewicht beizumessen. Ihr würde eine Parallaxe in Rektaszension von $1^{\mathrm{s}}5$ in Zeit entsprechen. Diese könnte man umso eher für vertretbar halten, als \textsc{Piazzi} aus den Beobachtungen in Palermo sogar das Doppelte gefunden hat. --- Der Stern $\delta$ in der Kassiopeia wurde bei 95 oberen und 54 unteren Kulminationen beobachtet, $\varepsilon$ im Großen Bären bei 26 oberen und 55 unteren Kulminationen. Die Beobachtungen fanden stets nur nachts oder in der Dämmerung statt. Daher hielt es \textsc{Oriani} hier nicht für nötig, Beobachtungen auszuschließen. Daraus ergab sich die Polhöhe
\begin{center}\small\begin{tabular}{@{}l@{\hspace{9pt}}r@{}}
aus $\delta$ in der Kassiopeia $.\;.\;.\;.\;.\;.\;.\;.\;.\;.$ & $45^\circ\,28'\,0''975$ \\
aus $\varepsilon$ im Großen Bären $.\;.\;.\;.\;.\;.\;.$ & $45\phantom{\circ}\,28\phantom{'}\,0.\,28$ \\
\end{tabular}\end{center}
\noindent und für 1811 ergaben sich die Deklinationen
\begin{center}\small\begin{tabular}{@{}l@{\hspace{9pt}}r@{}}
von $\delta$ in der Kassiopeia $.\;.\;.\;.\;.\;.\;.\;.\;.\;.$ & $59^\circ\,14'\,53''415$ \\
von $\varepsilon$ im Großen Bären $.\;.\;.\;.\;.\;.\;.$ & $56\phantom{\circ}\,59\phantom{'}\,15.\,19$ \\
\end{tabular}\end{center}
\noindent Genau genommen ist diese enge Übereinstimmung allerdings eher ein Beweis für die Güte der bei den Rechnungen verwendeten Refraktionstafel von \textsc{Carlini} als eine Bestätigung der Polhöhe. Wir würden uns wünschen, dass auch diese Tafel regelmäßig in den Ephemeriden erscheint. Zum Schluss fügt \textsc{Oriani} noch ein Verzeichnis der Deklinationen von 3c Zirkumpolarsternen bei (die drei oben genannten mitgerechnet), ohne jedoch die Ergebnisse der einzelnen Beobachtungen für sie anzugeben. Was die Polhöhe betrifft, stimmen, wie uns \textsc{Oriani} versichert, auch alle anderen Sterne bis auf einige Zehntel einer Sekunde mit dem oben angegebenen Ergebnis aus der Beobachtung des Polarsterns überein.

Der \emph{zweite} Aufsatz, ebenfalls von \textsc{Oriani}, bestimmt die Opposition des Mars im Jahr 1813. Dabei hatte \textsc{Carlini} die Rektaszensionen am Mittagsfernrohr, \textsc{Oriani} die Deklinationen am Kreis beobachtet. Als mittlerer Fehler der Marstafeln von \textsc{von Lindenau} ergaben sich geozentrisch $-4''8$ in der Länge und $+18''8$ in der Breite. --- Darauf folgt eine wertvolle Reihe von Meridianbeobachtungen der Sonne, die \textsc{Angelo Cesaris} von Anfang 1808 bis Ende 1811 am Mauerquadranten angestellt hat. Sie setzt die in den Jahrgängen 1809 und 1810 mitgeteilten Beobachtungen fort. --- Der anschließende Aufsatz von \textsc{Carlini} liefert Tafeln für die Mittelpunktsgleichung der Ceres und für die Reduktion dieses Planeten auf die Ekliptik. Die erste Tafel ist für die Exzentrizität $0.0784$, die zweite für die Neigung $10^\circ\,37'\,40''$ berechnet. Zugleich sind die entsprechenden Änderungen für eine Änderung der Exzentrizität um $0.0001$ und der Neigung um $10''$ angegeben. Ähnliche Tafeln für Pallas, Juno und Vesta werden für den nächsten Jahrgang angekündigt. Wir müssen gestehen, dass wir derartige Tafeln für Pallas und Juno nicht ganz zweckmäßig finden können, da die Theorie dieser Planeten doch auch in Zukunft die Form \emph{veränderlicher} Elemente beibehalten muss und die Veränderungen viel zu groß sind, als dass solche Tafeln die Rechnung erleichtern könnten. Um eine Vorstellung von der Größe dieser Veränderlichkeit zu geben, bemerken \SourceLeaf{581}wir, dass die Exzentrizität der Pallasbahn bei der Opposition von 1803 $=0.24554$, bei der Opposition von 1814 hingegen $=0.24135$ war. Bei der gegenwärtigen Art, die Planetenbewegungen zu berechnen, kann man auch auf diese Tafeln sehr gut verzichten. Die meteorologischen Beobachtungen an der Mailänder Sternwarte im Jahr 1813 beschließen den Band. Welch ein glückliches Klima für die praktische Astronomie, in dem man in einem einzigen Jahr 188 heitere Tage zählt!
